\color{unidarkgreen}{\subsection{Токовая защита нулевой последовательности стороны ВН (ТЗНП ВН)}}
\color{black}

\begin{enumerate}[label=\arabic{section}.\arabic{subsection}.\arabic*,  align=left, labelsep=4pt, leftmargin=0pt, itemindent=57pt]

\item
Токовая защита от замыканий на землю на стороне ВН используется в сетях с заземленной нейтралью
трансформатора. Защита предназначена для резервирования отключения замыкания на землю на шинах (и
линиях) стороны ВН трансформатора, а также для резервирования основных защит трансформатора.

\begin{figure}[H]
\centering
\includegraphics[width=1\textwidth,height=1\textheight,keepaspectratio]{\fbpath/021.0141 ТЗНП ВН/img/TZNP_VN.pdf}
\captionof{figure}{Функциональная схема ТЗНП ВН}\label{pict:TZNP_VN}
\end{figure}

\item
Защита подключается к ТТ, установленным со стороны обмоток ВН силового трансформатора.

\item
В состав ТЗНП ВН входит одна ненаправленная ступень с контролем тока нулевой последовательности
высокой стороны трансформатора. Порог срабатывания ступени защиты задается уставкой \textbf{«3I0ср ТЗНП ВН»},
а выдержка времени срабатывания ступени регулируется уставкой \textbf{«Тср ТЗНП ВН»}. Ввод в работу ТЗНП ВН
осуществляется с помощью программного переключателя \textbf{«ТЗНП ВН»}. Оперативный ввод/вывод ТЗНП ВН
производится с помощью БУ (блока управления функциями) «Вывод ТЗНП ВН». 

\item
Для обеспечения несрабатывания ТЗНП ВН от бросков намагничивающего тока предусмотрена
блокировка по содержанию второй гармоники в токе нулевой последовательности. Ввод в работу БНТ ТЗНП
ВН осуществляется с помощью уставки \textbf{«БНТ ТЗНП ВН»}. Блокировка основана на контроле отношения
действующего значения второй гармоники тока нулевой последовательности к основной гармонике. ТЗНП ВН
блокируется при превышении данного соотношения пороговой величины, заданной уставкой \textbf{«3I02г/3I01г
ТЗНП ВН»}. Для предотвращения излишней блокировки при больших токах КЗ, сопровождающихся
насыщением ТТ, блокировка выводится из действия при превышении тока 3I0 основной гармоники выше
уставки \textbf{«3I0ср БНТ ТЗНП ВН»}.

\item
Предусмотрена возможность оперативного ускорения ТЗНП ВН. Срабатывание оперативного ускорения
ТЗНП ВН действует на отключение трансформатора со всех сторон с пуском УРОВ и запретом АПВ. Ввод в
работу ОУ ТЗНП ВН осуществляется с помощью программного переключателя \textbf{«ОУ ТЗНП ВН»}, а выдержка
времени срабатывания ОУ ТЗНП ВН регулируется уставкой \textbf{«Тср ОУ ТЗНП ВН»}. Оперативный ввод/вывод ОУ
ТЗНП ВН производится с помощью команды управления \textit{«Ввод ОУ ТЗНП ВН»}.

\item
С целью исключения возникновения недопустимого режима работы параллельного трансформатора с
изолированной нейтралью в случае замыкания на землю одной фазы на шинах или участка сети стороны ВН,
резервная защита трансформатора с заземленной нейтралью выполняется действующей:
\begin{itemize}
\item с первой выдержкой времени на отключение (с пуском УРОВ и запретом АПВ) выключателей ВН трансформатора с изолированной нейтралью. Пуск ТЗНП ВН действует в логику отключения смежного Т, выдержка времени регулируется уставкой \textbf{«Тср откл см Т от ТЗНП»};
\item со второй выдержкой времени на разделение секции или систем шин ВН. Пуск ТЗНП ВН действует в логику деления стороны ВН, выдержка времени регулируется уставкой \textbf{«Тср ТЗНП откл СВ(ШСВ) ВН»};
\item с третей выдержкой времени на отключение (с пуском УРОВ и запретом АПВ) выключателей стороны ВН. Пуск ТЗНП ВН действует в логику деления стороны ВН, выдержка времени регулируется уставкой \textbf{«Тср ТЗНП откл В ВН»};
\item с выдержкой времени срабатывание ТЗНП ВН на отключение трансформатора со всех сторон с пуском УРОВ и запретом АПВ. Выдержка времени регулируется уставкой \textbf{«Тср ТЗНП ВН»}.
\end{itemize}

\item
В случае, когда оба трансформатора подстанции работают всегда с заземленными нейтралями, действие с
первой выдержкой времени на отключение параллельного трансформатора не используется.


\small
\setlength{\extrarowheight}{0.05cm} % высота ячеек в таблице
\begin{longtable}
{|>{\centering\arraybackslash}m{4.5cm}|>{\centering\arraybackslash}m{4.5cm}|>{\centering\arraybackslash}m{3cm}|>{\centering\arraybackslash}m{3.5cm}|}
\caption{Перечень уставочных параметров ТЗНП ВН\hfill\vspace{-0.5\baselineskip}}\label{tznp_vn:tbl1} \\ 
\hhline{|-|-|-|-|}
\rowcolor{gray!30}
Наименование уставки & Диапазон & Шаг &  Значение по умолчанию \\ 
\hhline{|-|-|-|-|}
\endfirsthead

\caption*{\hspace{3pt}\emph{Продолжение таблицы \ref{tznp_vn:tbl1}\hfill\vspace{-0.5\baselineskip}}} \\ % сделано по ГОСТ 2.105 п.6.8.7
\hhline{|-|-|-|-|}
\rowcolor{gray!30}
Наименование уставки & Диапазон & Шаг &  Значение по умолчанию \\  
\hhline{|-|-|-|-|}
\endhead
\hhline{|-|-|-|-|}

\endfoot
\endlastfoot

%%%%%%%%%%%%содержимое%%%%%%%%%%%%%%

\end{longtable}
\setlength{\extrarowheight}{0.0cm} % возврат параметра высота ячеек в таблице
\normalsize










































\end{enumerate} 